% ============================================================
%  Extended Summary (English) -- ~4-5 pages
%  To be translated to Greek and replace abstract_gr.tex
% ============================================================
\chapter*{Extended Summary}
\addcontentsline{toc}{chapter}{Extended Summary}

% ============================================================
\section*{Introduction and Motivation}

Point clouds have become the dominant representation for three-dimensional
geometric data across a broad range of applications: autonomous vehicles
equipped with LiDAR sensors, volumetric video capture systems for immersive
telepresence, and cultural heritage digitisation at millimetre precision.
In all of these settings, the sheer volume of raw data makes efficient
compression a prerequisite for practical use.

Traditional codecs developed by MPEG --- G-PCC for static geometry and V-PCC
for dynamic surfaces --- partially address this need but do not achieve the
rate-distortion efficiency that learned methods now routinely demonstrate on
image and video benchmarks.
Learned point cloud codecs, with PCGCv2~\cite{wang2021multiscale} as a
representative system, apply an analysis-synthesis paradigm directly to 3D
occupancy grids: an encoder maps the point cloud to a compact latent
representation, a learned entropy model estimates the latent distribution for
arithmetic coding, and a decoder inverts the process.
At high bitrates these systems match or exceed G-PCC.
At low bitrates, however, they reveal a consistent and under-studied failure
mode: \emph{geometric oversmoothing}.

Oversmoothing arises from a fundamental property of entropy models: rare,
fine-grained occupancy patterns incur high coding cost.
High-curvature regions --- edges, ridges, creases, corners --- produce exactly
such patterns in the occupancy grid.
At low bitrates, the encoder is implicitly penalised for representing these
regions faithfully, and the decoder reconstructs a surface in which such
features are blurred or erased.

The consequence is spatially non-uniform degradation: flat surfaces are
reconstructed near-perfectly while geometric detail concentrates the error.
Standard quality metrics such as point-to-point PSNR (D1)~\cite{tian2017geometric}
average over the entire reconstructed surface, weighting every point equally
regardless of its geometric role.
A codec that faithfully reconstructs 90\% of a surface while systematically
erasing all edges can still report a high aggregate PSNR --- the metric masks
precisely the information that matters perceptually.

% ============================================================
\section*{Contributions}

This thesis addresses two interrelated gaps: the \emph{inability to measure}
oversmoothing with existing metrics, and the \emph{absence of a correction method}
that operates on decoded point clouds without modifying the underlying codec.
The two principal contributions are:

\begin{enumerate}
  \item \textbf{Curvature-stratified PSNR metric.}
    An evaluation framework that partitions original point cloud points by local
    surface curvature and computes quality separately for edge and flat strata,
    exposing degradation that aggregate metrics cannot capture.

  \item \textbf{Rate-conditioned sparse U-Net post-processor.}
    A lightweight network that corrects geometric distortions in decoded point
    clouds without modifying the underlying codec, using a Feature-wise Linear
    Modulation (FiLM) head~\cite{perez2018film} to serve multiple operating
    rates with a single model.
\end{enumerate}

% ============================================================
\section*{Diagnostic Metric: Curvature-Stratified PSNR}

The metric operates in three steps: curvature estimation, stratification, and
per-stratum quality computation.

\paragraph{Curvature estimation.}
For each point $p$ in the \emph{original} cloud $\mathcal{P}$, local surface
curvature $\kappa(p)$ is estimated from the eigenvalue structure of the local
covariance matrix built from the $k=30$ nearest neighbours.
The smallest eigenvalue, normalised by the sum of all three, yields the
\emph{surface variation}: a value near 0 for planar regions and near $1/3$
for isotropically distributed neighbourhoods.

\paragraph{Stratification.}
Points are partitioned into flat, medium, and edge terciles using the 33rd
and 67th percentile curvature thresholds, computed per frame.
Percentile-based thresholds rather than fixed values make the stratification
dataset-adaptive: the absolute curvature range that constitutes an ``edge''
differs across subjects and scene types.

\paragraph{Quality direction.}
Quality is computed in the \emph{reverse} direction: from the original cloud
to the reconstructed cloud.
This measures how well the reconstruction \emph{covers} the original surface,
rather than whether reconstructed points are close to the original.
This direction is not merely a design preference but a necessity.
At low rates, the curvature of the \emph{decoded} cloud correlates negatively
(Pearson $r \approx -0.33$) with original curvature at edge regions, because
oversmoothing erases the fine structures that high curvature corresponds to.
Using decoded-cloud curvature for stratification would classify the
worst-affected points as flat, inverting the stratification at precisely
the rates where the metric matters most.

% ============================================================
\section*{Oversmoothing Characterisation in PCGCv2}

We apply the metric to PCGCv2 across four sequences of the 8i Voxelized Full
Bodies (8iVFB) dataset~\cite{8ivfb}: \textit{longdress}, \textit{loot},
\textit{redandblack}, and \textit{soldier}.
These sequences are the standard MPEG evaluation corpus for point cloud
compression and span a range of surface complexity.

The results reveal three systematic properties.

\paragraph{Systematic spatial bias.}
The flat-minus-edge PSNR gap
$\Delta_{\mathrm{smooth}} = \mathrm{PSNR}_{\mathrm{flat}} - \mathrm{PSNR}_{\mathrm{edge}}$
is strictly positive across all 28 sequence-rate combinations evaluated (seven
rate points, four sequences).
PCGCv2 \emph{always} reconstructs flat regions better than edge regions,
without exception.

\paragraph{The gap widens with rate.}
Counterintuitively, the oversmoothing gap \emph{increases} as the bitrate
increases.
At the highest evaluated rate (r7), the gap is 2--3$\times$ larger than at
the lowest (r1) for all sequences, ranging from 2.5 to 4.5~dB.
This demonstrates that oversmoothing is not simply a low-bitrate artifact
that disappears with additional bits: the codec allocates extra bits
preferentially toward flat surfaces, deepening the quality imbalance rather
than resolving it.

\paragraph{Sequence dependence.}
The gap magnitude varies substantially: \textit{redandblack} shows the largest
gap (2.4--4.5~dB), consistent with its fine surface detail (face, hair, fabric),
while \textit{loot} shows the smallest gap at low rates (0.2~dB at r1).

Applying the same evaluation pipeline to G-PCC reveals a qualitatively different
artifact profile: the gap peaks at intermediate rates (average 6.4~dB at ts4)
and then \emph{decreases} at higher rates, forming a non-monotone curve.
This contrast confirms that the monotonically growing gap is a structural
property of the learned occupancy-coding paradigm, not of geometry compression
in general.

\paragraph{Zero-inflation: implications for post-processing design.}
Examining the distribution of ground-truth displacement magnitudes
(nearest-neighbour vectors from decoded to original points) reveals that at
r2, 60\% of decoded points already coincide with original voxels and require
zero correction.
At r7, this fraction rises to 89\%.
This \emph{zero-inflation} property has a critical implication: a naive
regression model minimises its expected L1 loss by predicting zero for all
points.
The model collapses to a trivial solution that ignores the minority of points
where correction is most needed.

% ============================================================
\section*{Post-Processing Network}

\subsection*{Problem Formulation}

Let $\hat{\mathcal{P}} = \{\hat{p}_i\}_{i=1}^{N} \subset \mathbb{Z}^3$ be
the decoded point cloud at integer voxel coordinates (resolution $G=1024$)
and $\mathcal{P}^* = \{p^*_j\}_{j=1}^{M}$ the original reference cloud.
The post-processor predicts a 3D displacement $\Delta_i \in \mathbb{R}^3$
for each decoded point:
\[
  \hat{\mathcal{P}}_{\mathrm{refined}}
    = \bigl\{\hat{p}_i + \Delta_i\bigr\}_{i=1}^{N},
  \quad
  \Delta_i = f_\theta(\hat{\mathcal{P}},\,r)_i,
  \quad r = \log_2(\mathrm{bpp}).
\]
The output has the same number of points as the decoded cloud.
Displacements are bounded to ${\pm}5$ voxels via $\tanh$ activation, covering
all realistic corrections with headroom.
Input features are the three normalised voxel coordinates in $[-1,1]^3$;
an ablation study confirms that adding estimated curvature as an input feature
\emph{hurts} performance (discussed below).

\subsection*{Architecture: Three-Level Sparse U-Net}

The backbone is a three-level sparse U-Net operating directly on the decoded
occupancy grid, implemented with MinkowskiEngine~\cite{choy2019minkowski}.
Each encoder level halves the resolution and doubles the channel count
(32 $\to$ 64 $\to$ 128); the decoder mirrors this with skip connections.
The multi-scale structure provides both local geometric detail sensitivity
(edges) and broader surface context (topology), which a point-independent MLP
or global operation cannot achieve.

\subsection*{FiLM Rate Conditioning Head}

A two-layer MLP maps the scalar $r = \log_2(\mathrm{bpp})$ to 64-dimensional
FiLM parameters, split into 32 scale values $\boldsymbol{\gamma}$ and 32 shift
values $\boldsymbol{\beta}$, which modulate every active voxel's feature vector:
\[
  \mathbf{h}'_i = \boldsymbol{\gamma}(r) \odot \mathbf{h}_i + \boldsymbol{\beta}(r).
\]
The log-BPP encoding linearises the rate-distortion relationship so that equal
intervals correspond to approximately equal changes in distortion.
The FiLM MLP adds only 4,224 parameters (0.17\% of total model size),
enabling a single 2.5M-parameter model to serve all rate-distortion operating
points simultaneously.

\subsection*{Loss Function: Dynamic Chamfer Distance}

Standard displacement regression losses (L1, smooth-L1, curvature-weighted L1)
all fail due to the two structural problems identified above.

\emph{Zero-inflation failure}: at r2, the coordinate-wise median of ground-truth
targets is zero, so the optimal constant predictor under L1 is zero.
Any network trained with this loss collapses to near-zero predictions within
5--10 epochs.
This was verified empirically for all standard loss variants.

\emph{Nearest-neighbour noise}: at edge regions, a small perturbation of a
decoded point position can flip its nearest neighbour in the original cloud
from one side of the edge to the other, causing a discontinuous jump in the
target.
This instability is worst precisely where corrections are most needed.

We resolve both problems by abandoning precomputed per-point targets entirely.
Instead, we directly minimise the symmetric Chamfer Distance between the
refined patch and the corresponding crop of the original cloud:
\[
  \mathcal{L}_{\mathrm{CD}} =
    \frac{1}{|\hat{\mathcal{P}}_{\mathrm{ref}}|}
    \sum_{\hat{p} \in \hat{\mathcal{P}}_{\mathrm{ref}}}
      \min_{p^* \in \mathcal{P}^*_{\mathrm{crop}}} \|\hat{p} - p^*\|^2
    +
    \frac{1}{|\mathcal{P}^*_{\mathrm{crop}}|}
    \sum_{p^* \in \mathcal{P}^*_{\mathrm{crop}}}
      \min_{\hat{p} \in \hat{\mathcal{P}}_{\mathrm{ref}}} \|p^* - \hat{p}\|^2.
\]
The forward term penalises refined points far from the original surface;
the reverse term penalises surface regions not covered by the refined cloud.
The nearest-neighbour matching is recomputed at every training step on the
\emph{current refined positions} --- hence ``dynamic''.
For already-correct points the gradient contribution is near zero regardless
of the zero fraction, eliminating zero-inflation bias.
Nearest-neighbour noise is also resolved because the network can freely move
points within $[-5,5]^3$ and the matching adapts accordingly.

A critical implementation detail is that the original cloud crop must use
\emph{zero} boundary padding.
With non-zero padding, the reverse term is dominated by out-of-boundary
original points and collapses to near zero, causing a degenerate gradient that
pushes decoded points toward patch boundaries.
With zero padding, the two terms are balanced and the model learns correct
geometric corrections.

% ============================================================
\section*{Experimental Results}

\subsection*{Main Results on 8iVFB}

The final model is evaluated on 400 frames across all four 8iVFB sequences
at three rate points (r2, r4, r6).
Average PSNR improvements over the PCGCv2 baseline are:

\begin{center}
\begin{tabular}{lcc}
  \toprule
  Rate & Edge PSNR $\delta$ (dB) & Flat PSNR $\delta$ (dB) \\
  \midrule
  r2 ($\approx$0.05 bpp) & $\mathbf{+2.06}$ & $+1.47$ \\
  r4 ($\approx$0.15 bpp) & $+0.58$ & $+0.62$ \\
  r6 ($\approx$0.32 bpp) & $+0.27$ & $+0.18$ \\
  \bottomrule
\end{tabular}
\end{center}

At r2, the refined edge PSNR (62.11~dB) nearly closes the gap to the baseline
flat PSNR (62.70~dB).
The Bj{\o}ntegaard Delta PSNR integrated over r2--r6 is \textbf{+0.97~dB}
on the edge metric --- comparable in magnitude to two years of G-PCC
standardisation effort (approximately 1--2~dB BD-PSNR per version increment).
Improvements are consistent across all four sequences, including three not
seen during training (per-sequence range at r2: 1.78--2.21~dB edge improvement).

\subsection*{Ablation: Curvature as Input Feature}

Removing estimated curvature from the input features \emph{improves} results
at all rates.
At r2, edge PSNR improvement increases from +1.62~dB (with curvature) to
+2.06~dB (without), a gain of 0.44~dB.
Two factors explain this: (a) curvature estimated on the decoded cloud is
most unreliable at the high-curvature locations where corrections are most
needed; (b) the network's receptive field ($13^3$ voxels at full resolution)
already captures local surface geometry implicitly from the normalised
coordinate features.
This finding has a methodological implication: providing apparently relevant
features can hurt learning when those features are noisiest where the signal
matters most.

% ============================================================
\section*{Generalisation: Zero-Shot Transfer}

\subsection*{Unseen Codec: Unicorn}

The model is applied without retraining to Unicorn~\cite{wang2025unicorn},
a state-of-the-art learned codec not seen during training.
Results reveal a clear bitrate threshold.
At moderate rates (Unicorn R2--R4, bpp $>$ 0.05), distortion patterns differ
from PCGCv2's, causing slight edge degradation ($-$0.01 to $-$0.72~dB).
At extreme rates (R5--R6, bpp $<$ 0.025), oversmoothing is severe enough to
be codec-agnostic and the corrections transfer: \textbf{+1.09~dB} at R5 and
\textbf{+1.36~dB} at R6, averaged over all sequences.

\subsection*{Unseen Dataset: MVUB}

Evaluated on the MVUB dataset~\cite{loop2016mvub} (upper-body captures at
$G=512$, sequences \textit{david} and \textit{ricardo}), the model shows
positive edge improvement at low rates without retraining:
\textit{david} +0.92~dB and \textit{ricardo} +1.03~dB at r2.
At moderate-to-high rates the transfer is neutral to negative, mirroring
the Unicorn pattern.

Together, these experiments establish that low-rate geometric oversmoothing
is a cross-codec, cross-dataset failure mode, and that a model trained on a
single codec can generalise its corrections when the oversmoothing is severe
enough to dominate the distortion pattern.

% ============================================================
\section*{Conclusion}

This thesis addresses a specific and under-studied problem in learned point
cloud compression: geometric oversmoothing concentrated at high-curvature
surface regions, invisible to standard aggregate metrics.

We introduce a \emph{curvature-stratified PSNR metric} that makes this hidden
problem visible, and demonstrate that in PCGCv2 the oversmoothing gap is not
a low-rate artifact but a structural property of the learned occupancy-coding
paradigm that grows monotonically with rate.
The contrast with G-PCC --- which shows a non-monotone, rate-decreasing gap ---
confirms that oversmoothing is intrinsic to the occupancy-grid entropy model.

We propose a \emph{rate-conditioned sparse U-Net post-processor} with a FiLM
head, trained with a dynamic Chamfer Distance loss that circumvents the two
structural failure modes of standard displacement regression.
The model achieves substantial edge PSNR recovery at low rates (+2.06~dB at r2)
while also improving flat regions, and generalises without retraining to an
unseen codec and an unseen dataset at extreme compression rates.

\vspace{1em}
\noindent\textbf{Keywords:}
point cloud compression, post-processing, oversmoothing, curvature,
sparse convolutions, FiLM conditioning, rate-distortion, PCGCv2.
